% PDF-derived LaTeX transcription. Original manuscript source was unavailable.
% Published PDF SHA-256: 8aa114a924e374205d4029eb24c54e51248f4e2c1716c8e3da4226ac2b56ce8b
% Extracted text SHA-256: 84b47176f2eaa55d3ee4b542f3d72a797bcc4048335e3faa130eda9f17b80566
\documentclass[11pt]{article}
\usepackage[margin=1in]{geometry}
\usepackage{amsmath,amssymb}
\setlength{\emergencystretch}{3em}
\title{Independent Field and Charge Fidelity Versus Constructed Constraints in a Three-Dimensional Electromagnetic Spectral Codec}
\author{}
\date{}
\begin{document}
\maketitle
\noindent\textbf{PDF-derived LaTeX transcription.} This editable text follows the PDF extraction. Refer to the published PDF for figures, table layouts, and mathematical notation.

% Extracted chunk 1

Independent Field and Charge Fidelity Versus Constructed
Constraints in a Three-Dimensional Electromagnetic Spectral
Codec

\section*{Abstract}

Electromagnetic reconstruction can satisfy a calibrated energy envelope or a constructed
charge identity without recovering the reference field. We examine this distinction in an inher-
ited, parameter-free three-dimensional Yee-spectral codec using an explicitly selected set of 32
scenes representing 16 paired device cases. The design crosses four scene families, four inher-
ited split labels, and two propagation directions. Scene-level field normalized RMSE has mean
0.0306862 and empirical 95th percentile 0.0539417; complex spectral correlation has minimum
0.9981622, and weighted spatial phase RMS has maximum 0.0710899 rad. Nevertheless, target-
charge relative RMSE reaches 0.1427528 while the constructed decoded displacement-divergence-
minus-charge residual is zero. Reference-material energy retention lies between 0.9977533 and
0.9998071, but algebra shows that its near-unity upper envelope follows from a reference-derived
scale rather than independently demonstrating reconstruction or conservation. Raw receiver-
transfer reciprocity has maximum relative discrepancy 1.58973\ensuremath{\times}10\ensuremath{-}7 across 16 directional pairs.
A separate comparison of 563 and 643 raw transfers has maximum discrepancy 0.0616591; it
measures resolution sensitivity, not convergence order. Four archived analytic controls pro-
vide complementary, bounded evidence for wave structure, magnetic divergence, total internal
reflection, and interference. The contribution is an evidence-based separation of independent
fidelity, representation integrity, and raw-solver diagnostics, not a new codec or general physical
qualification.

\section*{1 Introduction and contributions}

A compressed electromagnetic field is useful only if the quantities needed for subsequent scientific
analysis survive reconstruction. Agreement with a scalar energy target is insufficient when the
scalar is itself used to normalize the reconstruction. Likewise, agreement between decoded charge
and decoded displacement divergence is not an independent charge-recovery test when charge is
computed from that displacement.
This study asks a specific question: what do the retained diagnostics of an inherited three-
dimensional spectral codec establish about independent field and charge fidelity, as distinct from
properties imposed by its representation or calibrated from its input? The object of study is
the inherited Yee-spectral codec with application binary interface version 4 (ABI4), which has no
learned parameters; we do not propose a new representation. Its temporal and spatial spectral
packets, material side channel, magnetic null-space coordinates, and finite coefficient packing are
inherited foundations.
Our contribution has three parts. First, we give the discrete and normalization conventions
needed to separate target-dependent fidelity from internal consistency. Second, we report recon-
1

struction measurements for a deliberately selected 4 \ensuremath{\times} 4 \ensuremath{\times} 2 design: 32 directional scenes and
16 paired device cases. Third, we interpret analytic controls, raw reciprocity, and two-resolution
transfer comparisons according to the quantities they actually measure. This separation matters
because strong internal consistency and measurable target error coexist in the same records.

\section*{2 Related work}

The inspected Meep author manuscript explains spatial and temporal staggering on the Yee grid
and the interpolation needed when combining field components or exporting fields [ 1]. That con-
text motivates retaining the inherited discrete operators here rather than substituting continuum
derivatives or treating all channels as collocated. Meep is not the solver evaluated in this study,
and its published properties do not validate the private implementation.
A companion retrospective manuscript, “Grid-native latent patchification for two-dimensional
Maxwell-field reconstruction: a checkpoint-selection-aware retrospective case study,” examines
learned two-dimensional convolutional representations, patch indexing, and deterministic correc-
tion. Here we instead examine an inherited three-dimensional parameter-free spectral codec and
distinguish reference fidelity from calibrated or constructed diagnostics. The available reports do
not provide a matched cross-codec comparison.
The pinned arXiv v1 treatment of waveguide reciprocity derives symmetry of a weighted scat-
tering matrix, CS = ( CS )T, under stated reciprocity assumptions [ 2]. Equality of ordinary scalar
transmission coefficients additionally depends on compatible modal and source–receiver normal-
ization. We therefore describe the measured forward–reverse scalar discrepancy as a raw-transfer
diagnostic, not a complete scattering-matrix test or a property of decoded fields.

\section*{3 Scientific methods and validation protocols}


\section*{3.1 Evidence and selected design}

This is a retrospective analysis of retained scalar diagnostics with source inspection; raw fields and
tokens were not retained for fresh replay, and no additional simulations were performed for this
analysis. Independent reductions of saved records check their arithmetic and identity joins, not the
original field measurements.
The numerical campaign contains 32 raw scene solves and 32 codec evaluations on 643 grids. It
selects one local directional pair per family and split combination. The families are homogeneous
wave, total-internal-reflection (TIR) interface, double slit, and waveguide scatterer. The inherited
split labels are train, validation, test, and ood (out-of-distribution); each family–split combina-
tion contains left\_to\_right and right\_to\_left scenes. Thus, there are 16 paired device cases,
not 32 independent devices. No model is trained on these split labels, and the labels do not make
the cases independent random validation samples.
The inherited raw source enables JAX 64-bit arithmetic and specifies Courant parameter 0.45,
1,024 native-grid time steps, and eight saved frames at steps 1017–1024. Its state consists of electric
displacement, magnetic field, and cell charge. Faraday uses forward staggered differences, while
Ampere and continuity use backward adjoints; conductivity is integrated by a midpoint rule. The
source declares periodic boundaries with an absorbing conductivity profile, rather than a CPML
claim. The separately recorded bridge environment used JAX 0.10.0 and Torch 2.10.0+cu128 on an
2

H100; its NumPy versions were 2.4.4 for production and 2.4.2 for consumption. These environment
observations are not a reconstruction of library consumption for every campaign record.

\section*{3.2 Inherited representation and units}

One logical scene has shape [8, 8, 64, 64, 64], with channels (Ex, Ey, Ez, Hx, Hy, Hz, \ensuremath{\rho}, \ensuremath{\varepsilon}r). Its native
packet has shape [4096, 256] and uses bfloat16 (BF16) storage. The supplied decoder source allo-
cates float32 field and material workspaces and complex64 periodic spectra; the magnetic null basis
is stored as complex64. These source declarations distinguish codec reconstruction precision from
the raw solver’s JAX 64-bit setting, but do not independently verify historical runtime precision.
Input and output use the same dimensionless code units; no SI conversion is supplied. Electric
components are scaled by 10\ensuremath{-}4, magnetic components by 2 \ensuremath{\times} 10\ensuremath{-}4, charge by 10\ensuremath{-}5, and material
values by (\ensuremath{\varepsilon}r \ensuremath{-} 1)/15.
The periodic packet uses an orthonormal spatial real FFT, a rank-two temporal basis, 55,000
adaptive field modes, and 60,000 independently selected charge modes. Signed coefficients have 255
levels in [\ensuremath{-}127, 127] and are packed two per finite BF16 word. Material geometry is stored at full
spatial resolution, but its values are not exact after reconstruction. Periodic magnetic coefficients
use two coordinates in the null space of the forward-divergence symbol, with an explicit remaining
DC coordinate. The separate open- x, periodic- yz lane forbids periodic projection. These are
representation choices, not new contributions of this study.

\section*{3.3 Discrete charge and magnetic diagnostics}

Let \ensuremath{\delta}+
i f (n) = f (n + ei) \ensuremath{-} f (n) and \ensuremath{\delta}\ensuremath{-}
i f (n) = f (n) \ensuremath{-} f (n \ensuremath{-} ei) denote the inherited code-unit
differences. For periodic boundaries these wrap. In the open- x lane, the terminal forward x
difference is zero and the initial backward x difference equals the initial value. Cell material is
averaged onto electric edges:
\ensuremath{\varepsilon}edge
i (n) = 1
2 [\ensuremath{\varepsilon}r(n) + \ensuremath{\varepsilon}r(n + ei)], (1)
G(E, \ensuremath{\varepsilon}r) =
X
i
\ensuremath{\delta}\ensuremath{-}
i (\ensuremath{\varepsilon}edge
i Ei), DH (H) =
X
i
\ensuremath{\delta}+
i Hi. (2)
The material shift at the terminal open- x edge repeats the terminal material value. There is no
implicit grid-spacing factor in G. A stored charge-operator multiplier \ensuremath{\alpha} \ensuremath{\in} \{ 1, 64\} supplies the
declared convention:
b\ensuremath{\rho} = \ensuremath{\alpha}G( bE, b\ensuremath{\varepsilon}r). (3)
Consequently, evaluating \ensuremath{\alpha}G( bE, b\ensuremath{\varepsilon}r) \ensuremath{-} b\ensuremath{\rho} with the same construction gives zero independently of
whether b\ensuremath{\rho} matches the supplied target \ensuremath{\rho}. Target-charge relative RMSE is
RMS(b\ensuremath{\rho} \ensuremath{-} \ensuremath{\rho})
max(RMS(\ensuremath{\rho}), 10\ensuremath{-}30) ,
with both RMS reductions taken over saved frames and spatial sites within a scene. A zero target
and zero reconstruction error therefore give zero through the denominator floor. The reported
values remain retained diagnostics, not freshly remeasured field comparisons.
For periodic synthesis, the stored charge spectrum first has its spatial mean removed. With
backward symbol pi = 1 \ensuremath{-} e\ensuremath{-}iki, the longitudinal correction is
\ensuremath{\Delta}Di(k) = p\ensuremath{\ast}
i \ensuremath{\Delta}\ensuremath{\rho}(k)P
j |pj|2 , k ̸= 0, (4)
3

Here \ensuremath{\Delta}\ensuremath{\rho} is the Fourier transform of the mean-removed stored desired charge divided by \ensuremath{\alpha}, minus
the current unmultiplied displacement divergence; the zero-mode correction is zero. The separately
recorded electric correction is the RMS change divided by the RMS pre-correction electric field,
with the denominator floored at 10\ensuremath{-}30. Both RMS reductions include all eight saved frames, three
electric components, and spatial sites within a scene. This synthesis precedes the per-frame energy
scaling in Section 3.5; final decoded charge is then derived from the scaled electric field and decoded
material using b\ensuremath{\rho} = \ensuremath{\alpha}G( bE, b\ensuremath{\varepsilon}r).
Magnetic divergence uses the forward symbol qi = eiki \ensuremath{-} 1. The periodic null basis satisfies
q · H = 0 without a post-decoding projection. The reported relative magnetic divergence is
rH = RMS(DH H)

P
i |\ensuremath{\delta}+
i Hi|2
1/2 , (5)
with denominator floored at 10\ensuremath{-}30. The open- x diagnostic excludes the first and last x planes.
Small rH tests this discrete constraint, not all Maxwell equations.

\section*{3.4 Independent field comparisons}

Write F = ( E, H ) and let S divide components by their fixed electric or magnetic scales. The
scene-level field metric is
eF = RMS[S( bF \ensuremath{-} F )]
RMS(SF ) . (6)
It aggregates the six field components, spatial sites, and eight saved frames within a scene, with
the denominator floored at 10\ensuremath{-}30. It is not an energy-weighted error or a pooled campaign voxel
error.
The retained spectral comparisons apply spatial real FFTs over the three grid axes, retaining
the saved-frame and field-component axes. The retained coefficients are flattened for scene-level
reductions, with no additional Hermitian multiplicity weighting. A reference-power mask excludes
coefficients below 10\ensuremath{-}10 of peak reference power. Writing Xk and bXk for reference and decoded
coefficients and wk = |Xk|2, the wrapped phase RMS is
" P
k\ensuremath{\in}M wk arg( bXkX \ensuremath{\ast}
k )2
max(P
k\ensuremath{\in}M wk, 10\ensuremath{-}30)
\#1/2
.
Complex spectral correlation is the magnitude of the masked complex inner product normalized by
the product of the masked spectral norms:

% Extracted chunk 4


P
k\ensuremath{\in}M X \ensuremath{\ast}
k bXk

% Extracted chunk 7


max
[U+0014][U+0000]P
k\ensuremath{\in}M |Xk|2[U+0001]1/2 [U+0010]P
k\ensuremath{\in}M | bXk|2
[U+0011]1/2
, 10\ensuremath{-}30
[U+0015] .
Degenerate reference spectra are rejected. These quantities expose phase and spatial structure that
a scalar energy multiplier cannot determine; their retained values were not freshly remeasured.

\section*{3.5 Reference-derived energy calibration}

The inherited per-frame energy is a material-weighted spatial mean in code units, with \ensuremath{\mu} = 1:
U (F ; \ensuremath{\varepsilon}) = 1
2
*X
i
\ensuremath{\varepsilon}edge
i E2
i +
X
i
H 2
i
+
. (7)
4

Let F0 be the reference and eF the reconstruction from the actual unscaled BF16 payload. Define
U0 = U (F0; \ensuremath{\varepsilon}), A = U ( eF ; \ensuremath{\varepsilon}), (8)
B = U ( eF ; b\ensuremath{\varepsilon}), s = c
p
U0/ max(A, B), c = 0.9999. (9)
The stored multiplier is a BF16 high part plus BF16 low residual, and bF = bs eF . For positive
reference energy, ignoring scale-storage rounding,
R\ensuremath{\varepsilon} = c2 A
max(A, B) , Rb\ensuremath{\varepsilon} = c2 B
max(A, B) , (10)
max(R\ensuremath{\varepsilon}, Rb\ensuremath{\varepsilon}) = c2 = 0.99980001. (11)
The smaller ratio depends only on agreement between the two material interpretations of the
unscaled field. Thus neither ratio independently constrains field shape or establishes a conservation
law. The implementation checks a retention lower bound of 0.995 and upper bound 1+2 \ensuremath{\times}10\ensuremath{-}6 after
BF16 storage. Zero reference energy has an exact-zero policy; invalid positive-energy reconstruction
fails rather than receiving an arbitrary scale.

\section*{3.6 Raw transfer and statistical summaries}

For a raw receiver trace an, the source computes
T = \ensuremath{\Delta}t
X
n
ane\ensuremath{-}i\ensuremath{\omega}(n+1)\ensuremath{\Delta}t, d (T1, T2) = |T1 \ensuremath{-} T2|
max(|T1|, |T2|, 10\ensuremath{-}30) . (12)
Reciprocity applies d to the two raw directions of each device case. Resolution sensitivity applies
it to matched 563 and 643 raw transfers. The latter report records 32 additional coarse solves and
zero additional native-grid solves, reusing retained native transfers. Neither comparison evaluates
decoded transfers.
Percentiles use linear interpolation between sorted scene or pair values: the fractional index
is (N \ensuremath{-} 1)q. They describe this explicitly selected set, including targeted difficult regimes; they
are not confidence intervals or population uncertainty estimates. Paired directions remain paired
observations.

\section*{4 Experimental results}


\section*{4.1 Selected-scene reconstruction}

Table 1 reports the retained scene-level summaries. High spectral correlation coexists with nonzero
target-charge error, despite the constructed zero Gauss residual; the latter is not an independent
measure of charge recovery. The largest field error, 0.0642646, and smallest spectral correlation
occur in the ood double-slit left-to-right case; its phase RMS is 0.0710899 rad. The opposite
direction has target-charge relative RMSE 0.1427528, the set maximum. All eight TIR records
report zero target-charge error, so the aggregate charge percentile includes zero-error cases rather
than describing only nonzero-charge scenes.
Material-side classification assigns each decoded permittivity to the nearest declared constituent
value. The misclassification fraction is the number of incorrect constituent labels divided by the
5

Table 1: Independent fidelity and material diagnostics for 32 selected scenes. Percentiles are
empirical set descriptions.
Metric Statistic Value
Field nRMSE Mean / p95 0.0306862 / 0.0539417
Complex spectral correlation Min. / p05 0.9981622 / 0.9986507
Phase RMS (rad) p95 / max. 0.0594879 / 0.0710899
Target-charge relative RMSE p95 / max. 0.1065627 / 0.1427528
Interface misclassification Max. fraction 0
Material absolute error Max. 0.0275002
Table 2: Raw-transfer discrepancies. Limits are the inherited p95 and maximum comparison thresh-
olds, not uncertainty intervals.
Comparison p95 Max. Limits
Reversal, 16 pairs 1.23979 \ensuremath{\times} 10\ensuremath{-}7 1.58973 \ensuremath{\times} 10\ensuremath{-}7 0.02 / 0.05
563 versus 643, 32 scenes 0.0541843 0.0616591 0.12 / 0.25
number of eligible reference cells, restricted to cells within 10\ensuremath{-}6 of a declared constituent value.
The reported zero fraction supports retention of constituent labels at these cells, but does not
establish exact material-value recovery or characterize mixed-material cells or subcell interface ge-
ometry. Similarly, the maximum periodic Gauss-synthesis electric correction, 0.0616582, is nonzero
even though the final constructed Gauss residual is zero. The maximum periodic magnetic input-
projection correction is 3.13419 \ensuremath{\times} 10\ensuremath{-}7, while maximum decoded relative magnetic divergence is
3.54912 \ensuremath{\times} 10\ensuremath{-}6.
Reference-material energy retention spans 0.9977533–0.9998071; decoded-material retention spans
0.9984918–0.9998074. These results are consistent with the declared calibrated envelope. They
must be read alongside, rather than substituted for, the independent errors in Table 1.

\section*{4.2 Raw reciprocity and resolution sensitivity}

Table 2 separates raw-solver transfer diagnostics from codec fidelity. Both reported summaries
are below their inherited comparison limits. The largest resolution discrepancy occurs in the ood
homogeneous-wave pair. Its two directions remain observations of the same device case. Agreement
under reversal does not make either resolution an independent reference solution.

\section*{4.3 Four analytic controls}

The archived CPU and GPU controls evaluate a periodic traveling wave, a discrete divergence-free
magnetic harmonic, an open- x TIR interface, and an open- x interference field. Table 3 gives repre-
sentative structure-sensitive results. These are four fixtures evaluated on two execution backends,
not eight independent physical trajectories.
For the traveling wave, the decoded temporal energy range relative to its mean is 5.82746 \ensuremath{\times} 10\ensuremath{-}7
on CPU and 3.49355 \ensuremath{\times} 10\ensuremath{-}7 on GPU. Its near-0.9998 total-energy retention remains calibrated ev-
idence; phase slope, envelope, and plane-flux comparisons provide complementary structure tests.
The magnetic harmonic retains RMS amplitude 0.9975800 on CPU and 0.9975799 on GPU, with
zero reported post-decoding projection correction. Perturbing all adaptive magnetic coefficient
6

Table 3: Archived analytic-control diagnostics. CPU and GPU values are separate retained obser-
vations, not error bars.
Control Quantity CPU GPU
Traveling wave Phase-slope relative error 2.01633 \ensuremath{\times} 10\ensuremath{-}8 1.94866 \ensuremath{\times} 10\ensuremath{-}8
Plane-flux mismatch 2.88467 \ensuremath{\times} 10\ensuremath{-}7 0
Magnetic harmonic Raw decoded divergence 1.17747 \ensuremath{\times} 10\ensuremath{-}7 1.16782 \ensuremath{\times} 10\ensuremath{-}7
Field correlation 0.9999855 0.9999855
TIR Fitted reflectance 0.9978352 0.9978352
Normal-power ratio 0.000252055 0.000252058
Interference Fringe amplitude retention 1.0000782 1.0000782
Bright-fringe phase RMS (rad) 0.000671750 0.000671739
values while forbidding the projection function still gives relative divergence 2.20142 \ensuremath{\times} 10\ensuremath{-}7 and
2.21850 \ensuremath{\times} 10\ensuremath{-}7, respectively. This supports the implemented null-coordinate property, not unre-
stricted electromagnetic validity. The magnetic fixture also records input-charge relative error 1.0
under its encoder diagnostic; its successful magnetic checks do not establish input-charge agreement.
The TIR control has zero interface transition width and zero side misclassification. Incident
and reflected amplitude retentions are approximately 1.000754 and 0.999670; evanescent-energy
retention is approximately 1.000499. Nonzero cladding field energy is allowed in TIR: the relevant
leakage diagnostic is normal time-averaged power, not erasure of the evanescent field. The fitted
reflectance differs from unity by approximately 0.00216477. The interference control preserves
intensity fringe bin 14 exactly and visibility retention approximately 0.99999984. The supplied
target harmonic residuals are finite, including magnetic-control Faraday relative RMS 6.39847 \ensuremath{\times}
10\ensuremath{-}5; the controls are discrete numerical fixtures, not exact continuum solutions.
Both backends also record exhaustive finite-BF16 round trips for all 65,025 signed-coefficient pairs
and successful shape, finiteness, and bit-transport checks. These establish bounded representation
behavior, not field accuracy by themselves.

\section*{4.4 Historical failure regimes and bridge scope}

The selectors explicitly include regimes associated with historical energy amplification 6.892087,
magnetic projection correction 0.481121, TIR interface width 31 cells, and fringe amplitude re-
tention 0.108828. Their selected real-source counterparts have field nRMSE 0.028564, 0.044933,
0.011154, and 0.064265, respectively. These are crosswalks to relevant regimes under a different
input contract, not before–after recovery of identical historical adversarial arrays.
The directly supplied standalone bridge report evaluates the double-slit counterpart and agrees
with its selected-set scalar record. A TIR scene in the selected set separately has field nRMSE
0.0101723, correlation 0.9999483, and phase RMS 0.00795672 rad. These observations demonstrate
bounded solver-to-codec evaluation of source-defined raw scenes; neither bridge evidence nor the
four controls expands the selected design into a comprehensive corpus evaluation.

\section*{5 Discussion and limitations}

The central result is not that constructed constraints are unhelpful, but that they answer different
questions. A finite, correctly decoded packet, an enforced magnetic null space, and a conservative
7

scale header are useful integrity properties. Independent field comparisons ask whether the reference
structure survives; target-charge error asks whether the supplied charge is recovered. The selected
records show why those questions cannot be collapsed: zero constructed Gauss residual coexists
with target-charge relative RMSE exceeding 0.14, and near-unity calibrated energy coexists with
nonzero field and phase errors.
The analytic controls add scientific specificity. Traveling-wave phase and flux comparisons test
more than amplitude normalization. The magnetic adversary tests the implemented divergence-
free coordinate construction without a hidden post-decoding repair. TIR distinguishes permitted
evanescence from normal power leakage, and the interference control tests fringe frequency, ampli-
tude, visibility, and phase. Their successes are meaningful within these fixtures; the finite target
residuals and input-charge discrepancy prevent treating the fixtures as universal physical certifi-
cates.
The retained evidence has material limits. Only one local pair is selected for each family–split
combination, including deliberately difficult selectors; these 16 device cases do not establish popu-
lation generalization or complete the 2,000-scene corpus. Only eight late-time frames are compared.
Raw reciprocity leaves decoded reciprocity undetermined. Two grid sizes quantify transfer sensitiv-
ity but cannot identify asymptotic convergence order or supply an independently verified reference
solution. Scalar retention limits independent remeasurement of spatial errors. Later reader re-
gressions support stricter stored-dtype and temporal-window semantics but do not retroactively
extend the numerical campaign to a different software candidate. Accordingly, the results support
a focused reconstruction characterization, not comprehensive electromagnetic qualification.

\section*{6 Conclusion}

For the selected 32-scene, 16-device design, the inherited spectral codec exhibits mean field nRMSE
0.0306862 and minimum complex spectral correlation 0.9981622, while target-charge relative RMSE
reaches 0.1427528. Those independent comparisons carry information that neither the reference-
calibrated energy envelope nor the constructed decoded charge identity can provide. Archived
analytic controls and raw-transfer diagnostics strengthen the bounded characterization when their
scopes are kept separate. The scientific contribution is a defensible interpretation of reconstruction
evidence, not a new codec, solver, conservation result, or certification.

\section*{7 Data and Code A vailability}

Supplementary record S1 is a private export containing archived source, numerical reports, source-
binding information, and independent scalar-checking recipes. The source includes the inherited
generator, codec, evaluator-related modules, and diagnostic scripts. The retained numerical arti-
facts are scalar reports; raw field and token arrays were not retained. The materials exist in the
private supplementary export, but no public repository, download link, or external access mecha-
nism has been established, and redistribution permission is not established. The recipes support
checking retained scalar arithmetic and source associations, not reconstructing the missing experi-
mental arrays.
8


\section*{8 References}

[1] A. F. Oskooi, D. Roundy, M. Ibanescu, P. Bermel, J. D. Joannopoulos, and S. G. Johnson.
Meep: A flexible free-software package for electromagnetic simulations by the FDTD method.
Computer Physics Communications 181, 687–702 (2010). Author manuscript dated 8 January
2010. https://doi.org/10.1016/j.cpc.2009.11.008.
[2] G. K. Svendsen, M. W. Haakestad, and J. Skaar. Reciprocity and the scattering matrix of waveg-
uide modes. arXiv:1301.2458v1 (2013 submission). https://arxiv.org/abs/1301.2458v1.
9

\end{document}
